\documentclass[12pt]{article}
\usepackage[T1]{fontenc}
\usepackage[utf8]{inputenc}
\usepackage[brazil]{babel}
\usepackage[margin=1in]{geometry}
\setlength{\parindent}{0pt}
\begin{document}
\section*{Tema 15}
Processo: PEDILEF 2008.72.95.001366-8/ SC\\
Situacao: Julgado\\
Ramo: DIREITO PREVIDENCIÁRIO\\
Questao: Saber se pode haver rateio de pensão entre esposa e concubina, no regime de concubinato impuro.\\
Tese: A pensão por morte não deve ser rateada entre a viúva e a concubina, pois a relação extraconjugal paralela ao casamento não configura união estável. Tese no mesmo sentido da firmada no Tema 526/STF.\\
Fonte: https://www.cjf.jus.br/phpdoc/virtus/
\end{document}
